\documentclass[12pt,a4paper]{article}
\usepackage[margin=1in]{geometry}
\usepackage{times}
\usepackage{parskip}
\usepackage{hyperref}
\usepackage{xcolor}
\usepackage{graphicx}

% Define colors
\definecolor{linkred}{RGB}{0, 45, 114}
\definecolor{linkblue}{RGB}{0, 45, 114}
\definecolor{jhuBlue}{RGB}{0, 45, 114}  % JHU Heritage Blue

\hypersetup{
    colorlinks=true,
    linkcolor=linkblue,
    urlcolor=linkblue
}

\pagestyle{empty}
\newcommand{\university}{University of Auckland}

\begin{document}


% Header with JHU logo
\begin{minipage}[c]{0.15\textwidth}
% Uncomment the line below and add your JHU logo file (download from brand.jhu.edu)
\includegraphics[width=2\textwidth]{jhu-icon.pdf}
\end{minipage}
\hfill
\begin{minipage}[c]{0.8\textwidth}
\raggedleft
{\Large \textbf{Jie Gao}}\\[0.3em]
{\color{jhuBlue}Malone Postdoctoral Fellow, Johns Hopkins University}\\
jgao77@jh.edu $|$ \href{https://gaojie058.github.io/}{gaojie058.github.io} $|$ +1 (410) 710-0555
\end{minipage}

\vspace{1.5em}
\noindent\rule{\textwidth}{0.4pt}
\vspace{1em}

January 2026

\vspace{1em}

Dear Members of the Search Committee,

I am writing to apply for the faculty position at \university. As a Malone Postdoctoral Fellow at Johns Hopkins University, working at the intersection of {\color{linkblue}Human-Computer Interaction},  {\color{linkblue}AI} and {\color{linkblue}Software Engineering}, I am deeply motivated by the challenge of designing AI systems that are not only powerful but also trustworthy and human-centred. My research vision, \textit{placing humans at the center of AI-supported work}, directly addresses one of the most pressing questions in computing today: How can we harness AI's efficiency while preserving human agency, interpretive judgment, and trust?

\textbf{Research Interests and Contributions}

My research focuses on {\color{linkred}human-AI collaboration} for unstructured data analysis. Over 90\% of organizational data is unstructured, including social media posts, clinical interviews, customer reviews, and code repositories, and its volume continues to grow rapidly. While AI offers efficiency, it often fails to capture the depth and nuance that human experts provide. My work addresses this gap by developing systems where AI acts as a collaborative partner, supporting human cognitive processes like sensemaking, reasoning, and reflection.

My foundational contributions include {\color{linkred}\textbf{CoAIcoder}} (TOCHI 2023, 100 citations), the first work to formally define AI-assisted human-to-human collaboration in qualitative data analysis, and {\color{linkred}\textbf{CollabCoder}} (CHI 2024, 181 citations), the first end-to-end LLM-powered workflow supporting all stages of collaborative qualitative analysis. I have also developed a systematic {\color{linkred}\textbf{taxonomy for human-LLM interaction modes}} (CHI EA 2024, 81 citations). My ongoing work includes {\color{linkred}\textbf{MindCoder.ai}}, a publicly accessible platform prioritizing verifiability over accuracy, and research extending human-AI collaboration to code comprehension (ICPC 2026).

\textbf{Future Research Agenda}

As a faculty member, I plan to lead a research lab advancing three interconnected directions: (1) {\color{linkblue}\textbf{Next Generation Human-AI Collaborative Systems}}, building theoretical and empirical foundations for human-AI joint activity and cognitive responsibility allocation; (2) {\color{linkblue}\textbf{Evaluation Frameworks for Human-AI Systems}}, developing rigorous frameworks combining qualitative research theory with scalable AI based evaluators; and (3) {\color{linkblue}\textbf{Broader Sensemaking Domains}}, extending text analysis to codebase understanding, AI agent trajectory analysis, and multimodal information interpretation.

\textbf{Interdisciplinary Experience and Collaboration}

I believe that achieving trustworthy human-AI collaboration requires a \textbf{multidisciplinary effort} that bridges AI, data, and system-centred fields (e.g., NLP, Software Engineering) with human-centred domains (e.g., Psychology, Social Sciences). My experiences across multiple institutions have equipped me with the mindset and skills needed to collaborate effectively across these disciplines. At {\color{jhuBlue}\textbf{Johns Hopkins University}}, I work with advisors spanning NLP (Mark Dredze), HCI (Ziang Xiao), and Robotics/Healthcare (Chien-Ming Huang). My postdoctoral work at the \textbf{Singapore-MIT Alliance (SMART)} with Thomas W. Malone provided foundational understanding of collective intelligence and cognitive science. I have collaborated with researchers across \textbf{5+ institutions} including visiting positions at NUS and Notre Dame, and co-organized the CHI 2024 workshop ``LLMs as Research Tools,'' fostering an interdisciplinary community around responsible AI usage in data work. In addition, my previous role at SMART, which also gave me first hand experience collaborating with industry partners such as \textbf{Singapore Changi Airport} and \textbf{Synergy Marine Group}, further equipped me with the ability to translate research into industry relevant problems.

\textbf{Alignment with \university}

My research aligns strongly with the {\color{linkblue}\textbf{School of Computer Science}} at the University of Auckland. The School's established HCI Research Group, with its focus on understanding and designing for human interaction with technology, resonates deeply with my approach of placing humans at the centre of AI-supported work. I am particularly excited about potential collaborations with Professor Paul Denny on computing education, particularly his research on conversations with Copilot and AI tools for scaffolding programming problem-solving. His flagship work on PeerWise, aligns well with my research philosophy of building real-world, impactful software tools to support learning and analysis. I also see opportunities to collaborate with Associate Professor Danielle Lottridge, whose work on human factors engineering and user behaviour aligns closely with my interest in how AI tools affect researcher cognition and analytical processes.

\textbf{Teaching and Mentoring}

I am committed to creating interactive, supportive learning environments. I have mentored 10 students across undergraduate, master's, and PhD levels, with several achieving significant outcomes including workshop publications and career advancement. I am prepared to teach core computer science and artifical intelligence courses, including \textbf{Introduction to Artificial Intelligence}, \textbf{Introduction to Software Engineering}, \textbf{Data Science}, \textbf{Human-Centered AI}, \textbf{Human-AI Interaction and Collaboration}. I am also able to provide other courses like User Research, Computational Qualitative Analysis, and Human Aspects of Software Engineering.



My long term goal is to develop new forms of critical human-AI collaboration that drive society's AI transformation by advancing efficiency, fostering trust, and generating social impact. I would welcome the opportunity to discuss how my research and teaching can contribute to \university's vision for human-centred computing and responsible AI.

\vspace{1em}

Sincerely,

Jie Gao

\end{document}